\odiseachapter{Decimotercera parte}{El asesinato de las niñas}

Recuerdo como si fuera ayer el día que terminé el canto del que nos vamos a ocupar en este capítulo. No pude seguir con el libro, asqueado, y no leería el canto XXIII ---el reencuentro con Penélope--- hasta muchos meses más tarde y lo hice sin ganas ni fe. Tenía quince años y ya entonces se la juré a Odiseo para siempre. Me costaría casi medio siglo ajustarle las cuentas en \emph{Abandonando Ítaca} y creo que todavía me quedé corto.

¿Por qué el asco, por qué el horror? En el capítulo anterior, Euriclea reconoce a Ulises, sangriento, después de la batalla. Este le pide a su nodriza que le informe sobre las mujeres de la casa:

\begin{verse}
De las mujeres de casa, ahora sí, dame cuenta cumplida,\\
de las que me hacen deshonra y las que de culpa están limpias.
\end{verse}

La nodriza acusa a doce esclavas y ofrece ir corriendo a despertar a Penélope, a la que hemos dejado durmiendo en sus altos aposentos.

\begin{verse}
Y el milardido Odiseo así respondiéndole dijo:\\
«No la despiertes aún; y di tú a las mujeres ya mismo\\
que vengan esas aquí, las que urdiendo vergüenza hayas visto».
\end{verse}

Lo que sigue es el episodio más espeluznante de toda la \emph{Odisea}. Odiseo le dice a su hijo y a sus dos fieles:

\begin{verse}
«Id ya sacando a los muertos, mandad a las criadas hacerlo,\\
y que después las trabancas y los sitiales rebellos\\
limpien frotando con agua y esponjas de mil agujeros.\\
Y cuando todo el casal en orden quede dispuesto,\\
sacad las esclavas del recio aposento y llevadlas al cerco\\
que hacen la espléndida valla del patio y el chozo de ruedo\\
y allí tajadlas a espada fileña hasta que el aliento\\
les arranquéis, y Afrodita se borre de sus pensamientos,\\
la que bajo los galanos en goce y a hurtadas tuvieron».
\end{verse}

Has leído bien, dolida lectora, herido lector. Odiseo instruye a sus compañeros para que, después de que las criadas hayan retirado los cadáveres y dejado todo bien limpio, las pasen a cuchillo. ¿Su pecado? Haberse acostado con los pretendientes. La crueldad no puede ser más extrema, el desprecio a la vida más alto, la sinrazón más absurda. No es de extrañar que Hollywood no se atreva ni a rozar semejante salvajada. Pero la \emph{Odisea} no puede entenderse sin el asesinato de las criadas.

Homero no nos ahorra ni un detalle:

\begin{verse}
Fue lo que dijo; y ya las mujeres en hato llegaban\\
gimiendo muy lastimeras, en llanto florido de lágrimas.\\
Allí los cuerpos recién fenecidos primero sacaban\\
al patio bienavallado y en el soportal los dejaban\\
unos sobre otros haciendo una pila: Odiseo orden daba\\
metiendo él mismo aguijón; los sacaban a fuerza obligadas.\\
Luego las sillas rebellas y las hermosas trabancas\\
con milojosas esponjas y agua purificaban.\\
Y el porquerizo y Telémaco y el guardador de las vacas\\
raspan a frote el solar de la bien hecha sala con palas,\\
y las esclavas la hez recogían y afuera la echaban.\\
Y cuando ya al fin dejaron en orden toda la estancia,\\
llevándose del aposento de sólido pie a las esclavas\\
al cerco del patio entre el chozo de ruedo y la espléndida valla,\\
en angostura las cierran, de donde ninguna escapara.
\end{verse}

Da igual que giman lastimeramente, da lo mismo que sean unas muchachas, casi unas adolescentes. Odiseo y sus hombres las tratan con menos respeto que a los cerdos de la piara de Eumeo, que a las vacas de Filetio, que a los perros del patio. Pero si te parece, angustiada lectora, ansioso lector, que la cosa no puede empeorar, siento desilusionarte. Ahora es Telémaco quien demuestra ser digno hijo de su padre:

\begin{verse}
Y es ahora el cauto Telémaco el que a los otros dos habla:\\
«No, no será en muerte limpia que yo les quite la vida\\
a quienes por mi cabeza vertieron tanta ignominia\\
y tanta en la de mi madre cuando con esos dormían».
\end{verse}

Es decir, nada de matarlas a espada como han hecho con los pretendientes. Después de todo, estos eran nobles y varones, es decir, se merecían una muerte limpia. En cambio, las pobres niñas, cuya insoportable afrenta se reduce a haber dormido con unos tipos que, posiblemente, tampoco les dieran mucha elección, se merecen algo peor. Las ahorca:

\begin{verse}
Eso decía, y ató de una nave proalazulosa\\
el cable a un alto pilar y tensándolo en torno a la choza,\\
bien en altura, dejoles sin suelo los pies allí a todas.\\
Como los tordos alidonosos o las palomas\\
que van a caer a una red armada por entre la fronda\\
cuando volvían al nido, y odioso lecho las toma,\\
ellas así sus cabezas en sarta tenían ---la gorja\\
de cada una en su nudo--- haciendo aún más triste su hora.\\
Y sacudieron sus pies una nada no más, vana y corta.
\end{verse}

La suerte del cabrero Melantio despeja toda duda sobre la clase de salvajes que son estos hombres:

\begin{verse}
Sacaron luego a Melantio por el umbral hasta el patio\\
y allí narices y orejas cortáronle en bárbaro tajo,\\
y las verijas le arrancan por crudas dar a los galgos,\\
y pies le truncan y manos con corazón enconado.
\end{verse}

El canto termina con la antítesis de las escenas horribles que acabamos de leer. Avisadas por su nodriza, llegan las esclavas fieles, para agasajar al héroe victorioso:

\begin{verse}
Y ya por todo el hermoso casal de Odiseo la anciana\\
marchó a avisar a las criadas para que allí se llegaran;\\
ellas, antorchas en mano, como en procesión por la sala,\\
en Odiseo se vierten enteras mientras lo abrazaban,\\
y su cabeza y sus hombros con mucha ternura besaban,\\
y acariciaban sus manos; y dulce lo asía a él la gana\\
de a lloro darse y sollozo: a todas llevaba en su alma.
\end{verse}

El golpe que me llevé de niño fue todavía mayor porque leí la \emph{Odisea} inmediatamente después de terminar la \emph{Ilíada}. Compara, agraviada lectora, ofendido lector, la conclusión de ambas épicas. Aquiles no es menos salvaje ni vengativo que Odiseo. El propio Héctor, cuyo nombre lleva mi hijo, no duda en matar a Patroclo, sin placer, pero sin piedad. El viejo Ulises que imaginé, reflexionando sobre su vida en Ítaca, recuerda también aquella tragedia:

\begin{verse}
Ese pobre muchacho, Patroclo,\\
mimado, dulce e inepto con las armas.\\
¿Creía que acostarse con el semidiós\\
bastaba para convertirlo en gran guerrero?\\
Con armadura y todo, Héctor lo redujo,\\
como un experimentado leñador\\
poda una planta joven. Ahora que lo piensas\\
era lo mismo cada día,\\
todos y cada uno, aqueos y troyanos,\\
muriendo a centenares, iban dejando atrás\\
un océano de lágrimas. Así son los caminos de la guerra.
\end{verse}

%El recuerdo de Pentesilea, la reina de las amazonas, se mezcla en su cansada memoria con el de Patroclo:
%
%\begin{verse}
%La espada del héroe había aniquilado a otros\\
%más dignos que el muchacho desafortunado.\\
%Pentesilea, por nombrar a una sola, la valiente amazona,\\
%que podía vencer a cualquier hombre ---tú incluido---,\\
%con la única excepción\\
%del Pelíada. Y sin embargo, como Héctor,\\
%luchó en una pelea que no podía ganar,\\
%eligiendo, igual que el príncipe troyano,\\
%el honor y la muerte.
%\end{verse}

Patroclo, cuyo único y fatal pecado fue creer que la magia de Aquiles le protegería:

\begin{verse}
¡Pobre Patroclo! Estaba convencido\\
de que la armadura y la reputación de Aquiles\\
lo protegerían. Y casi lo hicieron,\\
pero fue tan estúpido como para desafiar al troyano,\\
y pagó con su vida esa locura.
\end{verse}

En cambio, ni Héctor ni Pentesilea, la valiente amazona, se hicieron nunca ilusiones respecto al resultado de su encuentro con el Pélida:

\begin{verse}
En cambio, ni Héctor ni la reina de las amazonas\\
esperaban derrotar al guerrero invulnerable\\
y, sin embargo, se enfrentaron a él. ¿Por qué?
\end{verse}

¿Por qué, entonces, se pregunta el anciano Odiseo, se enfrentaron a él?

\begin{verse}
A menudo te has preguntado las razones,\\
no eres un cobarde, y sin embargo,\\
nunca habrías elegido una muerte segura,\\
cuando no era necesario. Aquiles no quería\\
matar a la valiente y hermosa mujer,\\
que en el campo de batalla era casi su igual,\\
no tenía con ella ningún problema, excepto el hecho aleatorio\\
de estar en bandos diferentes en una guerra que no era la suya.\\
Lloró sobre el cadáver,\\
casi tanto como cuando el soldado le trajo\\
el cuerpo de su amante.
\end{verse}

No, se dice el viejo Ulises, Aquiles no quería matar a Pentesilea. Tampoco a Héctor.

\begin{verse}
Comprender ese hecho tan simple te ha llevado décadas,\\
pero no has compartido con nadie tu conocimiento.\\
¿Para qué? Amigos y sirvientes por igual,\\
sacudirían sus cabezas y pensarían en un viejo que ha perdido el sentido.\\
Todo el mundo conoce la furia de Aquiles,\\
todo el mundo sabe cómo buscó a Héctor,\\
como un león busca una paloma,\\
o tal vez,\\
igual que un novio busca a su prometida. Tú estabas allí,\\
y oíste los gritos, y viste las lágrimas,\\
pero solo ahora entiendes que Aquiles\\
no lloraba a Patroclo,\\
sino contra los dioses, que le obligaron a matar\\
al único hombre de la Tierra\\
que merecía su amor.
\end{verse}

¿Amor? ¿Y por qué no? Ese distorsionado, inhumano, abominable amor que existe entre los héroes:

\begin{verse}
Quizás esa era la única razón por la que Héctor\\
le esperaba frente a las murallas,\\
en vez de ponerse a salvo. El amor,\\
ese distorsionado, inhumano, abominable amor,\\
que existe entre los héroes,\\
mientras intentan matarse el uno al otro.\\
También conoces ese amor fatal,\\
pues solo por milagro escapaste con vida\\
de la isla de Circe. Pero la verdad\\
es que no te habría importado morir en sus brazos,\\
como a manos de Aquiles a Héctor no le importó morir.
\end{verse}

\begin{verse}
Y murió. El hijo de Peleo lo mató, sin piedad,\\
y ató por los tobillos a su carro el cadáver\\
para dar así el más cruel espectáculo,\\
y de ese modo nadie sospechara\\
que tenía roto el corazón.
\end{verse}

\begin{verse}
Y entonces,\\
después de que las multitudes y los inmortales se alegrasen,\\
abrazó a Príamo y lloró con él,\\
antes de entregarle el cadáver de Héctor,\\
domador de caballos.
\end{verse}

Por supuesto, mi interpretación de los sentimientos de Aquiles ni es única ni tiene por qué ser la más acertada, pero lo que importa es que se basa en los mismos hechos que se describen en la \emph{Ilíada}. Héctor mata a Patroclo. Aquiles mata a Héctor y mancilla, furioso, el cadáver de su enemigo vencido, un acto que poco tiene que envidiar, en brutalidad, a las atrocidades del canto XXII de la \emph{Odisea}. Pero cuando Príamo pide compasión, Aquiles se apiada de él y le devuelve a su hijo para que pueda ofrecerle las honras fúnebres que se merece. Y esa nobleza reverbera para siempre en nuestra memoria colectiva. El héroe invencible, que en el campo de batalla segó más vidas que nadie, se vuelve humano con un solo acto de profunda compasión.

Naturalmente, Hollywood no tiene problema con Aquiles. La película \emph{Troya} es mucho menos fiel al texto de la \emph{Ilíada} que la de Nolan a la \emph{Odisea} y, sin embargo, la escena entre Príamo y Aquiles es memorable por la exactitud con que capta el espíritu de ambos hombres.

Sin embargo, es impensable que el blandicine de masas nos ofrezca una versión mínimamente fiel del canto XXII. La parte de la batalla no es el problema ---a pesar de lo cual Nolan no anda muy acertado con su Odiseo saltarín y su Telémaco pasmado---. El problema es mostrar al implacable asesino cortándole la cabeza a un suplicante o enfocarlo al final de la matanza cubierto de sangre, feroz como un león que acaba de devorar un toro. El problema es admitir la vileza con que Odiseo decide ejecutar a unas niñas. El problema es escamotearnos el detalle de que Telémaco ---al que solo le falta comer chicle en la Nodisea--- es un monstruo como su padre.

No es de extrañar que el cine tema a Homero, sobre todo al Homero de la \emph{Odisea}. Y lo paradójico es que, en realidad, teme a Odiseo no por ser un personaje primitivo, producto de una época remota, sino por su actualidad. Aquiles es un héroe a la antigua, que podemos reconocer, en buena parte, en Arturo o Lancelot. Su nobleza es admirable y del todo inverosímil. En cambio, la iniquidad de Odiseo es tan moderna como las guerras que asolaron el siglo XX y siguen asolando el XXI, como los campos de concentración, como las purgas de Stalin, como el Gran Salto Adelante. Odiseo nos recuerda demasiado a los ayatolás iraníes que también ahorcan mujeres por fruslerías.

No es de extrañar que el cine tema a Homero.
